% ============================================================================
%  SECCIÓN 1: INTRODUCCIÓN
% ============================================================================

\section{Introducci\'on}

La pr\'actica deportiva constituye una de las actividades de mayor impacto en la
salud, la convivencia social y el desarrollo humano de las comunidades. En este
contexto, la infraestructura deportiva \textemdash{}estadios, coliseos, canchas y
complejos polideportivos\textemdash{} cumple un rol central, pues representa el
espacio f\'isico donde se materializa la actividad deportiva y el encuentro de
deportistas, aficionados e instituciones. Sin embargo, la informaci\'on sobre
estos escenarios suele encontrarse dispersa, desactualizada o directamente
ausente, lo que impide que la ciudadan\'ia pueda conocer, localizar y acceder a
ellos de manera oportuna.

Los avances en tecnolog\'ias de la informaci\'on, en particular los sistemas de
informaci\'on geogr\'afica (SIG) y las aplicaciones m\'oviles, han demostrado ser
herramientas eficaces para organizar y difundir informaci\'on espacial de forma
intuitiva e interactiva \citep{longley2015}. Distintos sectores como el turismo,
el transporte y el comercio han transformado la manera en que las personas
acceden a los servicios mediante mapas interactivos y aplicaciones m\'oviles. No
obstante, el sector deportivo nacional a\'un presenta una brecha digital
significativa en este aspecto, sin que exista un medio centralizado que re\'una
la informaci\'on de los escenarios deportivos del pa\'is.

En respuesta a esta situaci\'on, el presente informe documenta la fase de
identificaci\'on de la problem\'atica y dise\~no inicial de la soluci\'on del
proyecto de desarrollo de un sistema web y m\'ovil para la visualizaci\'on
interactiva de escenarios deportivos a nivel nacional, denominado
provisionalmente \nombreApp{}. A lo largo del documento se describe la
problem\'atica identificada, la justificaci\'on del proyecto, el objetivo general
y los objetivos espec\'ificos, los usuarios, las funcionalidades preliminares, la
definici\'on del producto m\'inimo viable (MVP), los requerimientos funcionales y
no funcionales, el flujo de navegaci\'on, los wireframes iniciales y las
tecnolog\'ias propuestas para el desarrollo.

\subsection{Contexto general}

El deporte es reconocido como un derecho y un factor de desarrollo humano y
social, y su pr\'actica requiere de espacios adecuados, seguros y accesibles. A
nivel nacional existen m\'ultiples escenarios deportivos distribuidos en las
distintas ciudades y regiones del pa\'is, administrados por gobiernos
municipales, federaciones, clubes deportivos y entidades privadas. Estos
escenarios var\'ian en tama\~no, capacidad, disciplinas que albergan, estado de
conservaci\'on y servicios disponibles.

No se dispone, sin embargo, de un medio centralizado y accesible que re\'una
informaci\'on confiable sobre la ubicaci\'on, las instalaciones, los deportes
practicados, los horarios y los eventos de cada escenario. Esta situaci\'on afecta
tanto a los deportistas y aficionados, que desconocen la oferta deportiva de su
entorno, como a las propias entidades administradoras, que carecen de un canal
digital para dar visibilidad a su infraestructura y atraer a nuevos usuarios.

\subsection{Alcance del documento}

Este documento cubre la fase de identificaci\'on de la problem\'atica y dise\~no
inicial de la soluci\'on del proyecto. Est\'a dirigido al docente de la asignatura
y al equipo de desarrollo, y sirve como base para los siguientes avances del
proyecto. Describe la problem\'atica, los objetivos, los usuarios, las
funcionalidades preliminares, la definici\'on del MVP, los requerimientos
funcionales y no funcionales, el flujo de navegaci\'on, los wireframes
iniciales y las tecnolog\'ias propuestas. No incluye la implementaci\'on del
c\'odigo fuente, que corresponde a etapas posteriores del desarrollo.
